\documentclass[11pt]{article}
\usepackage[letterpaper,margin=1in]{geometry}
\usepackage[T1]{fontenc}
\usepackage[utf8]{inputenc}
\usepackage{amsmath,amssymb,amsthm,lmodern,microtype,booktabs}
\usepackage[hidelinks]{hyperref}

\newtheorem{theorem}{Theorem}
\newtheorem{lemma}[theorem]{Lemma}
\newtheorem{proposition}[theorem]{Proposition}

\title{Berge--Fulkerson Covers for the Mattiolo--Negrini--Pagani Snarks}
\author{}
\date{}

\begin{document}
\maketitle

\begin{abstract}
We prove that every snark $H_n$ ($n \ge 1$) in the cyclically $5$-edge-connected
unbounded-flow-resistance family of Mattiolo, Negrini, and Pagani~\cite{MNP} admits a
Berge--Fulkerson cover (six perfect matchings covering each edge twice). Specifically,
we verify computationally that the two finite building blocks $Z$ ($17$ vertices) and
$Y_1$ ($41$ vertices) each realize all $56$ admissible five-port boundary types, and we
give an explicit algebraic transition showing that the recursive layer preserves full
boundary realizability for every $Y_n$ ($n \ge 1$). Closing $Y_n$ into $H_n$ then yields
a Berge--Fulkerson cover for all $n \ge 1$.
\end{abstract}

\section{Introduction and Boundary Labels}

A \emph{Berge--Fulkerson coloring} of a cubic graph or multipole assigns to every edge
and semiedge a $2$-element subset of six colors $\{0,1,2,3,4,5\}$, encoded as a $6$-bit
incidence mask of weight~$2$, such that at every cubic vertex the three incident masks
partition the six colors. In a closed cubic graph, each color induces a perfect matching
and every edge belongs to exactly two of these six matchings.

Write $T = 63 = (111111)_2$ for the mask of all six colors and $\oplus$ for bitwise
exclusive or. Let $P$ be a five-pole (a cubic graph fragment with five dangling
semiedges at distinct terminals).

\begin{lemma}\label{lem:xor-admissible}
In any Berge--Fulkerson coloring of a five-pole $P$, the five boundary labels
$(\alpha,\beta,\gamma,\delta,\epsilon)$ satisfy
\[
  \alpha \oplus \beta \oplus \gamma \oplus \delta \oplus \epsilon = T.
\]
Conversely, any ordered $5$-tuple of weight-$2$ masks with exclusive or~$T$ corresponds,
up to permuting color names, to one of $56$ canonical boundary types. Consequently, a
five-pole realizing all $56$ boundary types accepts every ordered $5$-tuple of weight-$2$
masks with exclusive or~$T$.
\end{lemma}

\begin{proof}
Because $3|V(P)| = 2|E(P)| + 5$, the order $|V(P)|$ is odd. At each vertex $v \in V(P)$,
the three incident labels partition $T$, so their exclusive or is $T$. Taking the
exclusive or over all $v \in V(P)$ cancels every internal edge twice and leaves the
exclusive or of the five boundary labels on one side and $T$ (an odd number of copies) on
the other.

Conversely, if five weight-$2$ masks satisfy
$\alpha \oplus \beta \oplus \gamma \oplus \delta \oplus \epsilon = T$, then each color
appears at an odd number of ports (size $1$, $3$, or $5$) and the six port-support sizes
sum to $10$. Sorting the six port-support masks to forget color names leaves either the
unique all-five type $(5,1,1,1,1,1)$ or one of the $\binom{10+1}{2} = 55$ heavy types
$(3,3,1,1,1,1)$ determined by an unordered pair (with repetition) of the $10$ triples of
ports, with the four singleton multiplicities uniquely forced. Permuting the six colors in
a witness for the canonical type produces a coloring with the exact prescribed boundary
tuple.
\end{proof}

Our main result is the following.

\begin{theorem}\label{thm:main}
For every integer $n \ge 1$, the five-pole $Y_n$ of Mattiolo, Negrini, and
Pagani~\cite{MNP} realizes all $56$ five-port boundary types, and the closed snark $H_n$
(of order $40n + 2$) admits a Berge--Fulkerson cover.
\end{theorem}

\section{Finite Base Verification for \texorpdfstring{$Z$}{Z} and \texorpdfstring{$Y_1$}{Y1}}

We follow the construction of $M$, $N$, $Z$, $Y_1$, $Y_i$, and $H_n$ in Sections~3--4
and Figures~1--6 of~\cite{MNP}. Label the Petersen graph on $\{0,\dots,9\}$ with outer
$5$-cycle $0\text{--}1\text{--}2\text{--}3\text{--}4\text{--}0$, spokes $i\text{--}(i+5)$
for $0 \le i \le 4$, and inner $5$-cycle $5\text{--}7\text{--}9\text{--}6\text{--}8\text{--}5$
(encoded as \texttt{PETERSEN\_EDGES} in \texttt{verify\_flow\_resistance.py}):
\begin{itemize}
  \item $M$ (left of Figure~1 in~\cite{MNP}, $7$ vertices): delete the path $0,1,2$ and
        map $(e_1,\dots,e_5) \mapsto (4,5,3,7,6)$. The top port pair $e_1, e_3$ is
        adjacent, as is the bottom port pair $e_2, e_4$.
  \item $N$ (right of Figure~1 in~\cite{MNP}, $9$ vertices): delete vertex~$0$, delete
        edge $2\text{--}3$, and map $(f_1,\dots,f_5) \mapsto (2,3,1,5,4)$. The two
        trimmed-edge ends $f_1, f_2$ join $f_3, f_5$, respectively.
  \item $Z$ (Figure~2 in~\cite{MNP}, $17$ vertices): join $e_3$ to $f_1$ and $e_4$ to
        $f_2$, and add a cubic junction $j$ adjacent to $e_5$ and $f_5$, with boundary
        $(z_1,\dots,z_5) = (e_1, e_2, f_3, f_4, j)$.
  \item $Y_1$ (Figure~3 in~\cite{MNP}, $41$ vertices): combine a central $M$ with left
        and right copies $Z, Z'$ via joins $z_3 e_2$, $z_4 e_1$, $z'_3 e_4$, $z'_4 e_3$,
        and $z_5 z'_5$, with boundary
        $(\alpha,\beta,\gamma,\delta,\epsilon) = (z_1, z_2, z'_1, z'_2, e_5)$.
\end{itemize}

\begin{proposition}\label{prop:base}
Both $Z$ and $Y_1$ realize all $56$ five-port boundary types.
\end{proposition}

\begin{proof}[Computational verification]
For each pole and each odd port mask $S \subseteq \{1,\dots,5\}$, we remove the vertices
selected by $S$ and enumerate all perfect matchings of the remaining internal graph by
exact backtracking (minimum-degree branching), restoring the semiedges in $S$. For each
of the $56$ boundary types, an exact backtracking search selects one single-color
matching per color so that every internal edge is covered exactly twice, pruning whenever
an edge already covered twice is reused or remaining color domains cannot cover an edge
to multiplicity~$2$. Every returned labeling is independently checked for edge-label
validity, boundary agreement, and a partition of the six colors at every vertex. The
verification counts are summarized in Table~\ref{tab:base}.
\end{proof}

\begin{table}[ht]
\centering
\caption{Exact finite base verification for $Z$ and $Y_1$.}
\label{tab:base}
\begin{tabular}{@{}lrrr@{}}
\toprule
Pole & Vertices & Verified boundary types & Max single-color matchings (any fixed mask) \\
\midrule
$Z$   & $17$ & $56$ & $6$   \\
$Y_1$ & $41$ & $56$ & $109$ \\
\bottomrule
\end{tabular}
\end{table}

The standalone script \texttt{verify\_flow\_resistance.py} (requiring only the Python~3.10+
standard library) reproduces Proposition~\ref{prop:base}, verifies the transition
construction across all $56$ boundary types, and constructs and checks complete
six-matching covers for $H_1,\dots,H_8$ (orders $42$ through $322$, archived in
\texttt{certificates/bf-flow-resistance.json}):
\begin{verbatim}
python3 verify_flow_resistance.py --n 8
\end{verbatim}
Adding \texttt{--emit} outputs the full edge labels for all $56$ base covers of $Z$ and
$Y_1$ together with the six perfect matchings of each closed graph $H_1,\dots,H_n$.

\section{Uniform Extension Through the Recursive Layer}

In Figure~4 of~\cite{MNP}, the five-pole $Y_i$ ($i \ge 2$, order $40i + 1$) is formed
from $Y_{i-1}$, two new copies $Z$ (left) and $Z'$ (right) of $Z$, and six new cubic
junctions $A, B, C, D, E, F$ whose incidences are listed in Table~\ref{tab:junctions}.
The outer boundary of $Y_i$ is $(z_1, z_2, z'_1, z'_2, E)$, where the first four ports
belong to $Z$ and $Z'$ and the fifth port is the outer semiedge at~$E$.

\begin{table}[ht]
\centering
\caption{Incidences of the six junctions in the recursive layer from $Y_{i-1}$ to $Y_i$.}
\label{tab:junctions}
\begin{tabular}{@{}ll@{}}
\toprule
Junction & Incident neighbors or semiedges \\
\midrule
$A$ (top left)      & $z_3$, old $\gamma$, $C$ \\
$C$ (top right)     & $z'_3$, old $\delta$, $A$ \\
$B$ (bottom left)   & $z_4$, old $\alpha$, $D$ \\
$D$ (bottom middle) & old $\beta$, $B$, $E$ \\
$E$ (bottom right)  & $z'_4$, $D$, outer $\epsilon$ \\
$F$ (lowest)        & $z_5$, $z'_5$, old $\epsilon$ \\
\bottomrule
\end{tabular}
\end{table}

\begin{lemma}\label{lem:layer}
Let $(\alpha,\beta,\gamma,\delta,\epsilon)$ be any $5$-tuple of weight-$2$ masks with
$\alpha \oplus \beta \oplus \gamma \oplus \delta \oplus \epsilon = T$. Then there exist
weight-$2$ labels on the internal edges of the recursive layer and on the five old ports
of $Y_{i-1}$ such that:
\begin{enumerate}
  \item every junction $A, B, C, D, E, F$ has three incident masks partitioning $T$;
  \item the left copy $Z$ receives an admissible boundary $(\alpha,\beta,a,b,f)$ with
        exclusive or $T$;
  \item the right copy $Z'$ receives an admissible boundary $(\gamma,\delta,c,d,g)$ with
        exclusive or $T$; and
  \item the old boundary $(\alpha',\beta',\gamma',\delta',\epsilon')$ of $Y_{i-1}$ has
        exclusive or $T$.
\end{enumerate}
\end{lemma}

\begin{proof}
Observe first that every even-weight $6$-bit mask $M$ of weight at most~$4$ can be
written as $p \oplus q$ for two weight-$2$ masks $p, q$: if $|M| = 0$, take $p = q$ to be
any pair; if $|M| = 2$, split its two colors between $p$ and $q$ and add a common third
color to both; if $|M| = 4$, partition its four colors into $p$ and $q$.

We assign the labels in four steps:
\begin{enumerate}
  \item Choose a pair $d$ disjoint from $\epsilon$ and distinct from
        $\gamma \oplus \delta$. Since $\epsilon$ has weight~$2$, there are
        $\binom{4}{2} = 6$ pairs disjoint from $\epsilon$, and at most one equals
        $\gamma \oplus \delta$, so at least five valid choices exist. The even mask
        $M_R = T \oplus \gamma \oplus \delta \oplus d$ cannot equal $T$ (since
        $d \neq \gamma \oplus \delta$), so $|M_R| \le 4$. Choose pairs $c, g$ with
        $c \oplus g = M_R$; then $\gamma \oplus \delta \oplus c \oplus d \oplus g = T$.
  \item Similarly, choose a pair $f$ disjoint from $g$ and distinct from
        $\alpha \oplus \beta$ (at least five choices). Then
        $M_L = T \oplus \alpha \oplus \beta \oplus f$ has even weight at most~$4$, so
        there exist pairs $a, b$ with $a \oplus b = M_L$, making
        $\alpha \oplus \beta \oplus a \oplus b \oplus f = T$.
  \item Give edge $DE$ the mask $q = T \oplus d \oplus \epsilon$, which has weight~$2$
        because $d$ and $\epsilon$ are disjoint. Choose any pair $r$ disjoint from
        $a \cup c$ for edge $AC$, and any pair $s$ disjoint from $b \cup q$ for edge
        $BD$; both exist because $|a \cup c| \le 4$ and $|b \cup q| \le 4$, leaving at
        least two colors outside each union.
  \item At each of the six junctions, two of the three incident edges are now assigned
        disjoint pairs, forcing the remaining incident edge (an old port of $Y_{i-1}$) to
        take their complement in $T$:
        \[
          (\alpha', \beta', \gamma', \delta', \epsilon')
          = (T \oplus b \oplus s,\; T \oplus s \oplus q,\;
             T \oplus a \oplus r,\; T \oplus c \oplus r,\; T \oplus f \oplus g).
        \]
\end{enumerate}
Every entry of this old boundary is a weight-$2$ mask by the disjointness choices
($b \cap s = \varnothing$, $s \cap q = \varnothing$, $a \cap r = \varnothing$,
$c \cap r = \varnothing$, $f \cap g = \varnothing$), and every junction $A,\dots,F$
partitions $T$. Finally, canceling repeated terms $s$, $r$, and four copies of $T$, and
substituting $q = T \oplus d \oplus \epsilon$, we compute the exclusive or of the old
boundary:
\begin{align*}
  \alpha' \oplus \beta' \oplus \gamma' \oplus \delta' \oplus \epsilon'
  &= T \oplus a \oplus b \oplus c \oplus q \oplus f \oplus g \\
  &= (a \oplus b \oplus f) \oplus (c \oplus d \oplus g) \oplus \epsilon \\
  &= (T \oplus \alpha \oplus \beta) \oplus (T \oplus \gamma \oplus \delta) \oplus \epsilon \\
  &= \alpha \oplus \beta \oplus \gamma \oplus \delta \oplus \epsilon = T. \qedhere
\end{align*}
\end{proof}

\section{Proof of Theorem~\texorpdfstring{\ref{thm:main}}{1}}

\begin{proof}[Proof of Theorem~\ref{thm:main}]
We prove by induction on $n \ge 1$ that $Y_n$ realizes all $56$ boundary types. The base
case $n = 1$ is Proposition~\ref{prop:base}. For $n \ge 2$, fix any ordered $5$-tuple
$(\alpha,\beta,\gamma,\delta,\epsilon)$ of weight-$2$ masks with exclusive or~$T$. By
Lemma~\ref{lem:layer}, we obtain junction labels, admissible boundaries for both new
copies $Z$ and $Z'$, and an admissible boundary for $Y_{n-1}$. By
Proposition~\ref{prop:base} and Lemma~\ref{lem:xor-admissible}, both $Z$ and $Z'$ admit
Berge--Fulkerson colorings matching their prescribed boundaries; by the induction
hypothesis and Lemma~\ref{lem:xor-admissible}, $Y_{n-1}$ admits a coloring matching its
prescribed boundary. Gluing these colorings along the six junctions produces a
Berge--Fulkerson coloring of $Y_n$ with outer boundary
$(\alpha,\beta,\gamma,\delta,\epsilon)$. Thus $Y_n$ realizes all $56$ boundary types.

Finally, in Figure~6 of~\cite{MNP}, the closed snark $H_n$ is formed from $Y_n$ by
joining ports $\alpha$ and $\gamma$ into a single edge and attaching ports
$\beta, \delta, \epsilon$ to a new cubic vertex~$x$. Assign the same pair $\{0,1\}$ to
$\alpha$ and $\gamma$, and partition the six colors across $\beta, \delta, \epsilon$ as
$\{0,1\}, \{2,3\}, \{4,5\}$. The exclusive or of these five boundary masks is
$\{0,1\} \oplus \{0,1\} \oplus T = T$, so $Y_n$ admits a coloring with these boundary
labels. Gluing $\alpha$ to $\gamma$ gives an edge with matching label $\{0,1\}$, and the
three edges incident with $x$ partition the six colors. Hence $H_n$ admits a
Berge--Fulkerson cover for every $n \ge 1$.
\end{proof}

\begin{thebibliography}{9}
\bibitem{MNP}
D.~Mattiolo, A.~Negrini, and G.~Pagani,
\href{https://arxiv.org/abs/2604.22501}{arXiv:2604.22501}.
\end{thebibliography}

\end{document}
